% Complete proof module: Wilson construction and retained-coordinate comparison.
\section{A cyclic Wilson Leech lattice in the original triality space}
\label{wcl:sec:construction}

\subsection{The two retained metrics and the octonion convention}

Let $V=\mathbb O_1\oplus\mathbb O_2\oplus\mathbb O_3$ be the original
ordered octonion space. Retain its metric and cubic
\[
\begin{aligned}
 h_0(x,x)&=N(x_1)+N(x_2)+N(x_3),\\
 \tau(x_1,x_2,x_3)&=\operatorname{Re}((x_1x_2)x_3),\qquad
 d(Q_{\mathbb O}(0;x))=2\tau(x_1,x_2,x_3),
\end{aligned}
\]
where
\[
 Q_{\mathbb O}(\lambda;x)=
 \begin{pmatrix}\lambda&x_3&\bar x_2\\
 \bar x_3&\lambda&x_1\\x_2&\bar x_1&\lambda\end{pmatrix}.
\]
Wilson's norm on the same underlying ordered space is
\[
 q_{\rm W}(x)=\tfrac12 h_0(x,x).
\]
The map
\begin{equation}
 \iota:(V,q_{\rm W})\longrightarrow(V,h_0),\qquad
 \iota(x_1,x_2,x_3)=(x_1/\sqrt2,x_2/\sqrt2,x_3/\sqrt2)
 \label{wcl:eq:metric-isometry}
\end{equation}
is an explicit isometry. Both metrics remain specified. In particular,
for the original $\lambda$, which is not rescaled,
\begin{equation}
 d(Q_{\mathbb O}(\lambda;\iota x))
 =\lambda^3-\lambda q_{\rm W}(x)+\frac1{\sqrt2}\tau(x_1,x_2,x_3).
 \label{wcl:eq:cubic-scaling}
\end{equation}
Indeed the quadratic coordinates scale by $1/2$, while the trilinear
term scales by $(1/\sqrt2)^3=1/(2\sqrt2)$ and its original coefficient
is $2$. Thus the off-diagonal cubic after $\iota$ is
$1/(2\sqrt2)$ times its value before $\iota$.

For an explicit realization inside the same $\mathbb O$, choose a
Cayley--Dickson basis
$(1,\mathsf i,\mathsf j,\mathsf k,\ell,\mathsf i\ell,
\mathsf j\ell,\mathsf k\ell)$, with
\[
 (a+b\ell)(c+d\ell)=(ac-\bar d b)+(da+b\bar c)\ell,\qquad
 \mathsf i\mathsf j=\mathsf k .
\]
Wilson's ordered orthonormal basis is realized by
\begin{equation}
 (i_\infty,i_0,i_1,i_2,i_3,i_4,i_5,i_6)
 =(1,\mathsf i,\mathsf j,\ell,\mathsf k,
                  \mathsf j\ell,-\mathsf k\ell,\mathsf i\ell).
 \label{wcl:eq:actual-basis-map}
\end{equation}
The seven oriented Fano triples are $(t,t+1,t+3)$ modulo $7$:
$i_ti_{t+1}=i_{t+3}$, with cyclic products within each triple and
opposite signs on reversal. To verify \eqref{wcl:eq:actual-basis-map},
its seven corresponding products, in order $t=0,\ldots,6$, are
\[
 \mathsf i\mathsf j=\mathsf k,\quad
 \mathsf j\ell=\mathsf j\ell,\quad
 \ell\mathsf k=-\mathsf k\ell,\quad
 \mathsf k(\mathsf j\ell)=\mathsf i\ell,\quad
 (\mathsf j\ell)(-\mathsf k\ell)=\mathsf i,\quad
 (-\mathsf k\ell)(\mathsf i\ell)=\mathsf j,\quad
 (\mathsf i\ell)\mathsf i=\ell.
\]
Alternativity supplies cyclic products on each quaternionic line.
The seven lines cover all pairs of different imaginary basis units;
squares are $-1$ and $1$ is the unit. This proves the algebra
isomorphism on the full multiplication table, hence by bilinearity on
all octonions. It also preserves conjugation and the norm. No
conjugation of a coordinate of $Q_{\mathbb O}$ is performed.

\subsection{The explicit eight-dimensional lattice and two sublattices}

In the displayed Wilson basis write $e_0=i_\infty,e_{t+1}=i_t$ and put
\[
\begin{aligned}
 D_8&=\{(z_0,\ldots,z_7)\in\mathbb Z^8:\textstyle\sum z_i\equiv0\pmod2\},\\
 s&=\tfrac12(-1+i_0+i_1+i_2+i_3+i_4+i_5+i_6),\qquad
 L_{\rm W}=D_8+\mathbb Zs.
\end{aligned}
\]
This is Wilson's lattice $L$, renamed $L_{\rm W}$ to distinguish it
from the project's original rank-24 Niemeier lattice. Since $2s\in
D_8$ and $s\notin D_8$, it is exactly $D_8\cup(s+D_8)$.

\begin{lemma}\label{wcl:lem:e8}
The lattice $L_{\rm W}$ is even and unimodular, has minimum norm $2$,
and has the explicit basis
\[
 b_1=s,\quad b_j=e_{j-1}-e_j\ (2\le j\le7),\quad b_8=e_6+e_7.
\]
Its Gram matrix for the unchanged octonion norm is
\begin{equation}
 G_8=
 \begin{pmatrix}
 2&0&0&0&0&0&0&1\\
 0&2&-1&0&0&0&0&0\\
 0&-1&2&-1&0&0&0&0\\
 0&0&-1&2&-1&0&0&0\\
 0&0&0&-1&2&-1&0&0\\
 0&0&0&0&-1&2&-1&-1\\
 0&0&0&0&0&-1&2&0\\
 1&0&0&0&0&-1&0&2
 \end{pmatrix},\qquad \det G_8=1 .
 \label{wcl:eq:e8-gram}
\end{equation}
Its norm-two vectors are precisely the $112$ vectors $\pm e_i\pm e_j$
with $i<j$ and the $128$ vectors
$\frac12(\pm e_0\pm\cdots\pm e_7)$ with an odd number of minus signs.
\end{lemma}

\begin{proof}
For $z\in D_8$, $N(z)\equiv\sum z_i=0\pmod2$ and
$h(s,z)=\frac12\sum z_i-z_0\in\mathbb Z$.
As $N(s)=2$, every norm in $D_8+\mathbb Zs$ is even.
The covolume of $D_8$ is $2$, and its displayed extension has index
$2$, so $L_{\rm W}$ has covolume $1$. The integral bilinear form follows
by polarizing the even norm, and covolume one proves self-duality.
The columns $b_i$ all belong to $L_{\rm W}$. Their determinant has
absolute value one: expansion along the first coordinate gives
$1/2$ times the determinant $2$ of the indicated $D_7$ basis.
They therefore form a basis. Their scalar products give the displayed
matrix. Positivity and evenness give minimum at least $2$, and
$N(b_i)=2$ attains it. Integer vectors of norm two have exactly two
nonzero entries $\pm1$. Half-integer vectors have norm at least
$8/4=2$, with equality exactly when all eight entries are $\pm1/2$.
Their membership in $s+D_8$ is exactly the stated odd-minus parity.
\end{proof}

Define, with the bars in the positions printed in Wilson,
\[
 \mathsf M=L_{\rm W}\bar s,\qquad \mathsf N=L_{\rm W}s .
\]
All products are right products. Let $B$ be the basis matrix with
columns $b_i$. Right multiplication by $s$ in this basis is the
following integer matrix:
\begin{equation}
 S=B^{-1}R_s B=
 \begin{pmatrix}
 3&0&0&0&0&0&0&2\\
 -2&-1&0&1&-1&1&0&-2\\
 -4&-1&-1&1&0&0&1&-3\\
 -6&-1&-1&0&0&1&0&-4\\
 -8&0&-1&0&-1&1&1&-5\\
 -10&-1&0&0&-1&0&1&-5\\
 -5&0&-1&1&-1&0&0&-2\\
 -7&0&0&0&0&0&0&-4
 \end{pmatrix}.
 \label{wcl:eq:S}
\end{equation}
This is a direct multiplication-table calculation using
\eqref{wcl:eq:actual-basis-map}; it can be checked by multiplying each
column $b_j$ by the displayed $s$. Set $T=-I_8-S$. Since
$s+\bar s=-1$, $s\bar s=2$, and alternativity allows repeated $s$
and $\bar s= -1-s$, one has
\begin{equation}
 T=B^{-1}R_{\bar s}B,\qquad S+T=-I_8,\qquad ST=TS=2I_8,
 \qquad S^{\mathsf T}G_8S=T^{\mathsf T}G_8T=2G_8 .
 \label{wcl:eq:ST-identities}
\end{equation}

\begin{lemma}\label{wcl:lem:complementary}
Both $\mathsf M,\mathsf N$ are sublattices of $L_{\rm W}$ of index
$16$ and minimum norm $4$, and
\[
 2L_{\rm W}\subset\mathsf M,\mathsf N\subset L_{\rm W},\qquad
 \mathsf M+\mathsf N=L_{\rm W},\qquad
 \mathsf M\cap\mathsf N=2L_{\rm W}.
\]
Their images $\overline{\mathsf M},\overline{\mathsf N}$ in
$\mathcal E=L_{\rm W}/2L_{\rm W}$ are complementary
four-dimensional $\mathbb F_2$ subspaces.
\end{lemma}

\begin{proof}
The integer matrices $S,T$ give the inclusions in $L_{\rm W}$;
$ST=2I$ gives the inclusions of $2L_{\rm W}$.
Norm composition shows that each right multiplication scales all
squared lengths by $N(s)=N(\bar s)=2$; thus its determinant has
absolute value $(\sqrt2)^8=16$, proving the indices and minimum norms.
The identity $S+T=-I$ gives $\mathsf M+\mathsf N=L_{\rm W}$.
Modulo $2L_{\rm W}$, each image has size $2^8/16=2^4$.
Their sum is the eight-dimensional space $\mathcal E$, so their
intersection is zero. This proves the asserted intersection upstairs.
\end{proof}

\subsection{The binary glue code and its exact generators}

Wilson's original lattice is exactly
\begin{equation}
\begin{split}
 \Lambda_{\rm W}=\{(x,y,z)\in L_{\rm W}^3:\ &
 x+y,x+z,y+z\in\mathsf M,\quad x+y+z\in\mathsf N\}.
\end{split}
 \label{wcl:eq:Wilson-definition}
\end{equation}
It carries Wilson's norm $q_{\rm W}$, as printed. Define the lattice
inside the original triality metric by
\[
 \Lambda_{\rm cyc}=\iota(\Lambda_{\rm W})\subset(V,h_0).
\]
The simultaneous use of $\mathsf M=L_{\rm W}\bar s$ for pair sums
and $\mathsf N=L_{\rm W}s$ for the total sum is essential.

Write any class in $\mathcal E$ uniquely as $m+n$, with
$m\in\overline{\mathsf M}$ and $n\in\overline{\mathsf N}$.
Then the image of \eqref{wcl:eq:Wilson-definition} in $\mathcal E^3$
is the explicit binary code
\begin{equation}
 \mathcal C=
 \{(m_1+n,m_2+n,m_1+m_2+n):
       m_1,m_2\in\overline{\mathsf M},\
       n\in\overline{\mathsf N}\}.
 \label{wcl:eq:code}
\end{equation}
Indeed the pair conditions say that the three $\mathsf N$ components
are equal; the total-sum condition says that the sum of the three
$\mathsf M$ components is zero. This proves both directions of the
code description. Its dimension is $4+4+4=12$.

For a concrete lift, the first four columns of $S$ and of $T$ are
bases of their images modulo $2$. The determinant of the eight-column
matrix $(T_1,T_2,T_3,T_4,S_1,S_2,S_3,S_4)$ is $-7$, so these columns
are independent modulo $2$; each of the two images already has
dimension four. Put $m_j=b_j\bar s$, $n_j=b_js$ for $1\le j\le4$.
A complete set of $36$ integer generators of $\Lambda_{\rm W}$ is
\begin{equation}
\begin{gathered}
 (2b_i,0,0),\ (0,2b_i,0),\ (0,0,2b_i)\quad(1\le i\le8),\\
 (m_j,m_j,0),\ (m_j,0,m_j),\ (n_j,n_j,n_j)\quad(1\le j\le4).
\end{gathered}
 \label{wcl:eq:generators}
\end{equation}
These vectors are all in \eqref{wcl:eq:Wilson-definition}.
Their reductions span precisely \eqref{wcl:eq:code}; the first
$24$ generate its full kernel $(2L_{\rm W})^3$. This proves that
they generate the entire lattice, not only a sublattice.

\begin{theorem}\label{wcl:thm:Leech}
The explicitly generated lattice $\Lambda_{\rm W}$ with $q_{\rm W}$,
and therefore $\Lambda_{\rm cyc}$ with the unaltered $h_0$, is an
even unimodular rank-24 lattice of minimum norm $4$. In particular
it is rootless. This is the actual Wilson Leech construction, through
the explicit isometry \eqref{wcl:eq:metric-isometry}.
\end{theorem}

\begin{proof}
The code proves
$[L_{\rm W}^3:\Lambda_{\rm W}]=2^{24-12}=2^{12}$.
Since $L_{\rm W}$ has covolume one for its original octonion metric,
$\Lambda_{\rm W}$ has covolume $2^{12}$ for $h_0$.
Changing to the separately recorded metric $q_{\rm W}$ multiplies
volume in dimension $24$ by $(1/\sqrt2)^{24}=2^{-12}$,
so its covolume is exactly one.

For every triple of vectors in a real inner product space,
\[
 N(x)+N(y)+N(z)
 =N(x+y)+N(x+z)+N(y+z)-N(x+y+z).
\]
For a vector of $\Lambda_{\rm W}$ the four terms on the right
are norms in $\mathsf M$ or $\mathsf N$, so each is divisible by
$4$: these are twice norms in the even lattice $L_{\rm W}$.
It follows that $q_{\rm W}(x,y,z)$ is even. Polarization of these
even norms gives an integral bilinear form. An integral lattice of
covolume one equals its dual, proving unimodularity.

If a nonzero vector had $q_{\rm W}=2$, its three coordinate norms
would sum to $4$. Each nonzero coordinate in $L_{\rm W}$ has norm
at least $2$, so one coordinate is zero. The other two then lie in
$\mathsf M$, each of whose nonzero vectors has norm at least $4$.
At most one can be nonzero. That coordinate also lies in $\mathsf N$
by the total-sum condition, so it lies in
$\mathsf M\cap\mathsf N=2L_{\rm W}$ and has norm at least $8$,
a contradiction. Hence all nonzero norms are at least $4$.
The vectors $(2b_i,0,0)$ have $q_{\rm W}=4$, so the minimum is
exactly $4$. The isometry $\iota$ transfers every statement to the
original metric.
\end{proof}

\subsection{Cyclic symmetry, the fixed lattice, and exact three-primary glue}

Retain the canonical cycle $c(x,y,z)=(z,x,y)$. Denote the opposite
cycle used in Wilson coordinates by $c_{\rm W}(x,y,z)=(y,z,x)$.
Direct substitution gives
\[
 c_{\rm W}=c^{-1}=c^2,\qquad c_{\rm W}^2=c,\qquad
 \langle c_{\rm W}\rangle=\langle c\rangle.
\]
Both actions are retained with these directions; no action is renamed
by silently reversing it.
The defining three coordinate memberships, three pair conditions,
and total-sum condition are all unchanged by $c_{\rm W}$; hence both cycles preserve
$\Lambda_{\rm W}$ and $\Lambda_{\rm cyc}$. They preserve their metrics
and the original cubic:
\[
 \operatorname{Re}((xy)z)=\operatorname{Re}((yz)x).
\]
This follows from the octonion adjoint identities
$h(ab,d)=h(b,\bar a d)=h(a,d\bar b)$, which imply real associativity,
and $\operatorname{Re}(ab)=\operatorname{Re}(ba)$.
In the Albert algebra $c_{\rm W}$ is realized by conjugation with the
\emph{real} permutation matrix for $(1\,3\,2)$.
It sends the coordinates $(x_1,x_2,x_3)$ to $(x_2,x_3,x_1)$ with no
conjugations, and permutes the ordered diagonal frame accordingly.
The inverse real permutation matrix, for $(1\,2\,3)$, realizes the
canonical $c$ and sends $(x_1,x_2,x_3)$ to $(x_3,x_1,x_2)$.
Thus the fixed lattices of the two cycles are identical. For the
retained residue isometry $\mathcal Lr=c\mathcal L$, the opposite
action satisfies $\mathcal Lr^{-1}=c_{\rm W}\mathcal L$.

\begin{theorem}\label{wcl:thm:fixed-coinvariant}
Before applying $\iota$, the cyclic fixed lattice and its orthogonal
complement in $\Lambda_{\rm W}$ are exactly
\[
 F=\{(n,n,n):n\in\mathsf N\},\qquad
 C=\{(m_1,m_2,-m_1-m_2):m_1,m_2\in\mathsf M\}.
\]
Their ranks are $8$ and $16$. For the retained Wilson metric, and
therefore after $\iota$ for the original triality metric, their Gram
matrices in the displayed bases are
\begin{equation}
 G_F=3G_8,\qquad
 G_C=\begin{pmatrix}2G_8&G_8\\G_8&2G_8\end{pmatrix}.
 \label{wcl:eq:fixed-coinvariant-grams}
\end{equation}
The fixed-lattice basis is $(b_js,b_js,b_js)$, $1\le j\le8$;
the coinvariant bases are $(b_j\bar s,0,-b_j\bar s)$ and
$(0,b_j\bar s,-b_j\bar s)$. In particular the fixed lattice has
minimum norm $6$ and determinant $3^8$.
The coinvariant lattice has determinant $3^8$, and the map which
negates its second $L_{\rm W}$ parameter identifies its displayed
Gram matrix with the usual $A_2\otimes E_8$ Gram matrix.
\end{theorem}

\begin{proof}
A fixed triple has the form $(x,x,x)$. Its pair condition holds for
every $x\in L_{\rm W}$ because $2L_{\rm W}\subset\mathsf M$.
Its total condition is $3x\in\mathsf N$. As $L_{\rm W}/\mathsf N$
has exponent two, $3x\in\mathsf N$ is equivalent to $x\in\mathsf N$.
This proves the first description. Orthogonality to the whole
diagonal real eight-space says $x+y+z=0$. The pair conditions are
then exactly $x,y,z\in\mathsf M$, and the total condition is automatic.
This proves the description of $C$.

Right multiplication by either $s$ or $\bar s$ multiplies the
octonion metric by $2$. Thus
\[
 q_{\rm W}(as,as,as)=3N(a),
\]
and
\[
 q_{\rm W}(a\bar s,b\bar s,-(a+b)\bar s)
 =N(a)+N(b)+N(a+b)=2N(a)+2N(b)+2h(a,b).
\]
Polarization gives both full Gram matrices. Their determinants are
\[
 \det G_F=3^8\det G_8=3^8,\qquad
 \det G_C=\det\begin{pmatrix}2&1\\1&2\end{pmatrix}^{8}
 (\det G_8)^2=3^8.
\]
Minimum norm $6$ for $F$ follows from the exact
factor $3$ and the minimum $2$ in $L_{\rm W}$.
The sign map $b\mapsto-b$ changes both off-diagonal blocks from
$G_8$ to $-G_8$, proving the stated exact comparison to
$A_2\otimes E_8$.
\end{proof}

There is a complete integral comparison, including the glue:
\begin{equation}
 0\longrightarrow F\oplus C\longrightarrow\Lambda_{\rm W}
 \xrightarrow{\ \Sigma\ }\mathsf N/3\mathsf N
 \longrightarrow0,\qquad
 \Sigma(x,y,z)=x+y+z\pmod{3\mathsf N}.
 \label{wcl:eq:three-primary-exact}
\end{equation}
To prove surjectivity, for every $a\in L_{\rm W}$ the explicit vector
\[
 v(a)=(as,a,-a)
\]
lies in $\Lambda_{\rm W}$ and has sum $as$. Its pair sums are
$a(s+1)=-a\bar s\in\mathsf M$,
$a(s-1)=-a\bar s-2a\in\mathsf M$, and $0$;
its total lies in $\mathsf N$. If a lattice vector has total $3n$,
subtracting $(n,n,n)\in F$ leaves a vector of $C$. This proves the
kernel assertion. Thus
\[
 \Lambda_{\rm W}/(F\oplus C)\cong\mathsf N/3\mathsf N
 \cong L_{\rm W}/3L_{\rm W}\cong(\mathbb Z/3)^8,\qquad
 [\Lambda_{\rm W}:F\oplus C]=3^8.
\]
All maps commute with $c_{\rm W}$ and with its inverse $c$; both
cycles act trivially on the quotient.

The discriminant anti-isometry is also explicit. Use the quadratic
discriminant convention $q_D(u+D)=h_D(u,u)\pmod{2\mathbb Z}$
for an even lattice $D$ with its retained metric. The orthogonal
projections of the lift $v(a)$ are
\begin{align}
 f(a)&=\tfrac13(as,as,as),\nonumber\\
 \gamma(a)&=\left(\tfrac23as,\ a-\tfrac13as,\ -a-\tfrac13as\right).
 \label{wcl:eq:glue-projections}
\end{align}
They belong to $F^*$ and $C^*$, respectively: their pairings with
$F$ or $C$ are the pairings of the integral vector $v(a)$ with the
same lattice, because the other projection is orthogonal.
The classes $f(a)+F$ identify $L_{\rm W}/3L_{\rm W}$ with $F^*/F$.
If $\gamma(a)\in C$, then $f(a)=v(a)-\gamma(a)\in
\Lambda_{\rm W}\cap F_{\mathbb R}=F$, so $a\in3L_{\rm W}$.
Since $C^*/C$ has order $\det G_C=3^8$, the map
\[
 f(a)+F\longmapsto\gamma(a)+C
\]
is an isomorphism of the two discriminant groups. Its quadratic
sign is determined without an unspecified convention:
\[
 q_F(f(a))=\tfrac13N(a),\qquad
 q_C(\gamma(a))=q_{\rm W}(v(a))-\tfrac13N(a)
              =2N(a)-\tfrac13N(a).
\]
As $N(a)$ is even, their sum $2N(a)$ is zero modulo $2\mathbb Z$.
This is the required anti-isometry. Equations
\eqref{wcl:eq:three-primary-exact} and
\eqref{wcl:eq:glue-projections} specify the overlattice itself,
including its actual coset representatives.

\subsection{Finite lattice symmetries and exact cubic tests}

Independent sign changes of the three sectors preserve
$\Lambda_{\rm W}$. Changing the sign of $x$ changes a pair sum by
$-2x\in2L_{\rm W}\subset\mathsf M$ and changes the total by
$-2x\in2L_{\rm W}\subset\mathsf N$. The same argument applies to
each sector. All six sector permutations preserve the defining
conditions too. Under signs $(\epsilon_1,\epsilon_2,\epsilon_3)$
the cubic $\tau$ is multiplied by $\epsilon_1\epsilon_2\epsilon_3$;
hence exactly the even sign changes in this subgroup preserve it.
Cyclic permutations preserve $\tau$ as already proved.

The odd permutation $(x,y,z)\mapsto(y,x,z)$ need not preserve it,
even on the lattice itself. In Wilson's basis put
\[
 a=2(1+i_0),\qquad b=2(1+i_1),\qquad d=2(i_2+i_3).
\]
Every coordinate belongs to $2L_{\rm W}$, so $(a,b,d)\in
\Lambda_{\rm W}$. The real part of
$((1+i_0)(1+i_1))(i_2+i_3)$ is $-1$, because
$i_0i_1=i_3$ and the $i_3^2$ term is the only real contribution.
Reversing the first two factors changes this contribution to $+1$.
Consequently
\begin{equation}
 \tau(a,b,d)=-8,\qquad \tau(b,a,d)=8.
 \label{wcl:eq:transposition-cubic}
\end{equation}
After the specified scalar isometry the original determinant cubic
has values $-4\sqrt2$ and $4\sqrt2$.

The simultaneous Fano-index operations
\[
 A:i_t\longmapsto i_{t+1},\qquad
 B_0:i_t\longmapsto i_{2t},\qquad i_\infty\longmapsto i_\infty
\]
are octonion algebra automorphisms, as the defining oriented Fano
products show. They are orthogonal permutations preserving $D_8$ and
the odd-minus half coset, and they fix the displayed $s$ and $\bar s$.
They consequently preserve $\mathsf M,\mathsf N,\Lambda_{\rm W}$.
Applied simultaneously to all three sectors, they preserve $\tau$.
They satisfy $A^7=B_0^3=1$ and $B_0AB_0^{-1}=A^2$; their permutations
are the distinct maps $t\mapsto 2^j t+k$ with $j=0,1,2$ and
$k\in\mathbb F_7$. Thus they form an explicitly realized
$C_7\rtimes C_3$ of order $21$.

The even sign changes and the cycle $c_{\rm W}$ form
$(\mathbb Z/2)^2\rtimes C_3\cong A_4$. They commute with the
simultaneous Fano operations, and their intersection is trivial:
a simultaneous Fano operation fixes the real unit in every sector
and does not permute sectors. Hence an explicit finite group
\[
 A_4\times(C_7\rtimes C_3)
\]
of order $252$ preserves both the literal lattice and the original
cubic. This is a proved subgroup; no assertion that it is the full
cubic stabilizer is needed.

Wilson also gives the lattice symmetries
\[
 r_t(x,y,z)=(x,yi_t,zi_t)\qquad(t=0,\ldots,6).
\]
For completeness their lattice action can be checked solely from
the explicit finite matrices above. Let
$U_t=B^{-1}R_{i_t}B$. The signed permutation multiplication table
and \eqref{wcl:eq:S} give
\begin{equation}
 U_t\in GL_8(\mathbb Z),\qquad
 T^{-1}(U_t-I_8)\in M_8(\mathbb Z),\qquad
 \tfrac12(U_t-I_8)T\in M_8(\mathbb Z)
 \quad(0\le t\le6).
 \label{wcl:eq:unit-matrix-certificates}
\end{equation}
These are finite exact matrix assertions, not assumptions about
octonionic matrix associativity. The full multiplication rule and
the exact rational algorithm checking every entry for all seven
values are included in the accompanying \texttt{check\_wilson.py}.
Concretely, start with the $64$ products
$1i_t=i_t1=i_t$, $i_t^2=-1$, and the seven oriented Fano triples,
form the eight columns $R_{i_t}(e_j)$, conjugate by the explicitly
displayed $B$, and substitute the integer $T=-I_8-S$.
For each of the three matrices in \eqref{wcl:eq:unit-matrix-certificates}
the reduced denominators of every entry are $1$, for each
$t=0,1,2,3,4,5,6$. Also $\det U_t=1$.
This is a reproducible $7\cdot3\cdot64$-entry certificate.
The first condition is also immediate from the signed permutation:
right multiplication by $i_t$ flips four coordinate signs, preserving
the odd-minus parity of $L_{\rm W}$.

The second matrix assertion says
$L_{\rm W}(i_t-1)\subset\mathsf M$, and the third says
$\mathsf M(i_t-1)\subset2L_{\rm W}$. Thus $x+yi_t$ and $x+zi_t$
differ from the original pair sums by elements of $\mathsf M$.
The pair $(y+z)i_t$ remains in $\mathsf M$, and the transformed total
differs from $x+y+z$ by
$(y+z)(i_t-1)\in2L_{\rm W}\subset\mathsf N$.
All defining conditions are preserved. The inverse follows from
$r_t^2(x,y,z)=(x,-y,-z)$ and the proved sign symmetries.
Norm composition proves that $r_t$ is an isometry.

Nevertheless none of these seven $r_t$ preserves the original cubic.
For $a_t=2(1+i_t)$, the triple $(a_t,a_t,a_t)$ lies in
$(2L_{\rm W})^3\subset\Lambda_{\rm W}$. All products below occur in
the associative real subalgebra generated by $i_t$, and give
\[
 \tau(a_t,a_t,a_t)=8\operatorname{Re}((1+i_t)^3)=-16,
\]
whereas $(1+i_t)i_t=-1+i_t$ and
\[
 \tau(a_t,a_ti_t,a_ti_t)
 =8\operatorname{Re}((1+i_t)(-1+i_t)^2)=16.
\]
Therefore the original determinant cubic on the isometric lattice
$\Lambda_{\rm cyc}$ changes from $-8\sqrt2$ to $8\sqrt2$.
This supplies actual lattice vectors as counterexamples, not only
ambient real vectors outside the lattice.

\subsection{Exact comparison with coordinate lattices and the marked lane}

The construction has the explicit integral sandwich
\begin{equation}
 (\sqrt2L_{\rm W})^3
 \ \subset\ \Lambda_{\rm cyc}\
 \subset\ (L_{\rm W}/\sqrt2)^3,
 \qquad
 [\Lambda_{\rm cyc}:(\sqrt2L_{\rm W})^3]
 =[(L_{\rm W}/\sqrt2)^3:\Lambda_{\rm cyc}]=2^{12}.
 \label{wcl:eq:coordinate-sandwich}
\end{equation}
This is \eqref{wcl:eq:code} transported by the exact isometry
\eqref{wcl:eq:metric-isometry}. The inclusions are literal inclusions
in the fixed three octonion sectors, and each commutes with both
$c_{\rm W}$ and the canonical $c=c_{\rm W}^{-1}$.
Together with the explicit bases, the code, and the three-primary
glue \eqref{wcl:eq:glue-projections}, it determines actual
commensurability maps and cosets rather than only an abstract
isometry class.

Under \eqref{wcl:eq:actual-basis-map}, the single negative column
takes the Wilson odd-minus half roots to even-minus half roots in
the current ordered Cayley--Dickson basis. The element $s$ remains
exactly
\[
 s=\tfrac12(-1+\mathsf i+\mathsf j+\mathsf k+\ell
                     +\mathsf i\ell+\mathsf j\ell-\mathsf k\ell),
\]
with coordinate vector $(-1,1,1,1,1,1,1,-1)/2$.
The signed permutation has determinant $+1$ and all $64$ basis
products agree, as independently checked in the accompanying code.
Thus the lattice, its cyclic action, and the triality cubic live
in one specified multiplication convention.

\subsection{Literal intersections with both retained lattice embeddings}
\label{wcl:sec:literal-intersections}

Let $U=\operatorname{span}_{\mathbb R}(1,\mathsf i,\mathsf j,\mathsf k)$
be the first four coordinates in the \emph{current} ordered
Cayley--Dickson basis, and let
\[
 D_4=\{a=(a_0,a_1,a_2,a_3)\in\mathbb Z^4:
                          a_0+a_1+a_2+a_3\in2\mathbb Z\}\subset U.
\]
This quaternionic $D_4$ denotes the displayed four-coordinate lattice;
the original Niemeier $D$ summand below remains in its original place.

Let $J_{\rm W}$ denote the signed permutation matrix in
\eqref{wcl:eq:actual-basis-map}. If $U_0:\mathbb R^4\to\mathbb R^8$
is the coordinate inclusion in the current basis, then
\begin{equation}
 S^{-1}B^{-1}J_{\rm W}^{-1}U_0=
 \begin{pmatrix}
 1/2&-1/2&-1/2&-1/2\\
 -1/2&0&0&1/2\\
 -1&1/2&0&1/2\\
 -3/2&1&1/2&1/2\\
 -2&1&1&1/2\\
 -5/2&3/2&1&1\\
 -5/4&3/4&3/4&3/4\\
 -7/4&3/4&3/4&3/4
 \end{pmatrix}.
 \label{wcl:eq:quaternion-membership-matrix}
\end{equation}
All matrices on the left have already been explicitly specified.
Thus the equation is a direct rational multiplication certificate.

\begin{lemma}\label{wcl:lem:quaternion-intersection}
In this fixed coordinate frame,
\[
 \mathsf N\cap U=2D_4 .
\]
\end{lemma}
\begin{proof}
Since $\mathsf N\subset L_{\rm W}$, an element of $\mathsf N\cap U$
has integral first four coordinates: the half-integer coset of the
current $E_8$ description has no zero coordinate among its last four.
Write those integral coordinates as $k_0,k_1,k_2,k_3$. Membership in
$\mathsf N$ is exactly integrality of the matrix
\eqref{wcl:eq:quaternion-membership-matrix} applied to $k$.
Subtracting its last two rows gives $k_0/2$, so $k_0$ is even.
The second row then forces $k_3$ even, the third forces $k_1$ even,
and the fourth forces $k_2$ even. Write $k=2a$.
The first six rows are now integral. The seventh and eighth rows
are integral precisely when
$-5a_0+3a_1+3a_2+3a_3$ is even, equivalently
$a_0+a_1+a_2+a_3$ is even. Conversely these conditions make all eight
rows integral. This proves both inclusions with no unstated lattice
identification.
\end{proof}

We now retain the original lattices and their actual map $\mathcal L$.
The original coordinates are $X=(x_{i,j};d_k)$,
$1\le i\le4$, $0\le j\le5$, $1\le k\le4$, with
$\sum_jx_{i,j}=0$. The root lattice is
$R=A_5^{\perp4}\perp D$, with its unchanged Euclidean metric.
Its index-$72$ overlattice $N$ is the inverse image of the original
glue
\[
 (k_1,k_2,k_3,k_4;y)
   =\bigl(3u+2(a+b,a+2b,a,b);\,(u_2+u_4,u_3+u_4)\bigr),
\]
where $u\in\mathbb F_2^4$ has even weight and $a,b\in\mathbb F_3$;
the first four entries are modulo $6$. Here $k_i=-6x_{i,0}\pmod6$
and $y$ is the original $D^\vee/D$ marking, so $d=0$ has $y=00$.
Retain the original neighbour
\[
 \rho_A=(5,3,1,-1,-3,-5)/2,\quad
 \rho_D=(3,2,1,0),\quad \rho=\rho_A^{\oplus4}\perp\rho_D,
\]
\[
 K=\{X\in N:(\rho,X)\in6\mathbb Z\},\qquad
 \alpha=(e_0-e_1,0,0,0;0),\qquad
 \Lambda=K+\mathbb Z(\rho/6-\alpha).
\]
Thus $\rho\in N$, $\rho^2=84$, $(\rho,\alpha)=1$.
Every vector of $N$ or $\Lambda$ has rational original coordinates.

For exact use of the retained isometry, set
\[
 \sigma_i=\sum_{j=0}^5(-1)^jx_{i,j},\qquad
 t_{i,p,j}=x_{i,p+2j}-\frac13\sum_{h=0}^2x_{i,p+2h}.
\]
Write $u_0=1,u_1=\mathsf i,u_2=\mathsf j,u_3=\mathsf k$ and retain
the remaining current Cayley--Dickson basis vectors $u_4,\ldots,u_7$.
The already specified map has the exact coordinate form
\[
 (\mathcal LX)_{j+1,a}=
 \begin{cases}
 t_{i,p,j}+\sqrt2\,\sigma_{a+1}/6,&0\le a\le3,\\
 t_{i,p,j}+\sqrt3\,d_{a-3}/3,&4\le a\le7,
 \end{cases}
 \qquad a=2(i-1)+p.
\]
The placement of $\sigma_{a+1}$ is retained, including its difference
from the index $i$ of the residual. No coordinate is reassigned.

\begin{theorem}\label{wcl:thm:literal-intersection}
The literal intersections in the original metric space are exactly
\begin{equation}
 \Lambda_{\rm cyc}\cap\mathcal L(N)
 =\Lambda_{\rm cyc}\cap\mathcal L(\Lambda)
 =\mathcal I
 :=\{\sqrt2(a,a,a):a\in D_4\subset U\}.
 \label{wcl:eq:literal-common-lattice}
\end{equation}
In original Niemeier coordinates the inverse image of this common
lattice is
\[
 \mathcal L^{-1}(\mathcal I)
 =\{x_{i,j}=(-1)^j a_{i-1},\ d=0:\ a\in D_4\}\subset K.
\]
It has rank four, minimum norm $12$, and Gram matrix
\begin{equation}
 G_{\mathcal I}=6
 \begin{pmatrix}
 2&-1&0&0\\
 -1&2&-1&-1\\
 0&-1&2&0\\
 0&-1&0&2
 \end{pmatrix},
 \qquad\det G_{\mathcal I}=4\cdot6^4=5184.
 \label{wcl:eq:common-gram}
\end{equation}
An explicit basis is $\sqrt2(\delta_j,\delta_j,\delta_j)$,
where
$\delta_1=(1,-1,0,0)$, $\delta_2=(0,1,-1,0)$,
$\delta_3=(0,0,1,-1)$, $\delta_4=(0,0,1,1)$.
\end{theorem}

\begin{proof}
Each coordinate of $\Lambda_{\rm cyc}$ belongs to
$\mathbb Q/\sqrt2=\sqrt2\mathbb Q$ in the current frame.
For a rational original vector $X$, the first four coordinate
formulas above therefore force $t_{i,p,j}=0$ at the corresponding
positions, by linear independence of $1,\sqrt2$ over $\mathbb Q$.
The last four force both $t_{i,p,j}=0$ and $d_{a-3}=0$, by linear
independence of $1,\sqrt2,\sqrt3$. For the latter independence, a
relation with nonzero $\sqrt3$ coefficient would place $\sqrt3$ in
$\mathbb Q(\sqrt2)$; squaring $A+B\sqrt2$ forces $AB=0$ and then
would make $3$ or $3/2$ a rational square, a contradiction.
All residuals therefore vanish and $d=0$. The zero fibre sums give
\[
 x_{i,j}=(-1)^j\sigma_i/6.
\]

First let $X\in N$. Its $A_5^\vee$ integral-difference conditions
are exactly $k_i:=\sigma_i/3\in\mathbb Z$, yielding
$x_{i,j}=(-1)^jk_i/2$. Their discriminant labels are
$-3k_i=3k_i\pmod6$. Their three-primary parts are zero, so the
injectivity of $(a,b)\mapsto(a+b,a+2b,a,b)$ forces $a=b=0$.
Since $d=0$, the binary condition is
$u_2+u_4=u_3+u_4=0$ and $\sum u_i=0$. These equations give
$u_1=u_2=u_3=u_4$. Thus all four $k_i$ have the same parity.
Conversely these conditions satisfy every original glue requirement.
The image is the diagonal triple
\[
 \mathcal LX=(k/\sqrt2,k/\sqrt2,k/\sqrt2),\qquad k\in U.
\]
By the fixed-lattice description of Wilson's construction it
belongs to $\Lambda_{\rm cyc}$ exactly when $k\in\mathsf N$.
The previous lemma makes this equivalent to $k=2a$, $a\in D_4$,
which proves the first intersection equality.

For the second equality, a vector of $\Lambda$ can be written
$X=n+m\rho/6$ with $n\in N$, $m\in\mathbb Z$.
The already proved necessary condition $d_X=0$ gives
$d_n=-m(3,2,1,0)/6\in D^\vee$.
All coordinates of a vector of $D^\vee$ are integral or all
half-integral; the fourth coordinate is zero, so all are integral.
Its third coordinate forces $6\mid m$, and hence $X\in N$.
In fact $\Lambda\cap N=K$: in
$k+m(\rho/6-\alpha)\in N$ the pairing with the original root
$\alpha$ forces $6\mid m$, and
$(\rho,\rho-6\alpha)=84-6=78\in6\mathbb Z$.
Thus an intersection vector belongs to $K$.
For the diagonal-fibre vectors just found,
\[
 (\rho,X)=\frac32\sum_{i=1}^4 k_i.
\]
Consequently the condition $X\in K$ is $\sum k_i\equiv0\pmod4$.
Every $k=2a$ with $a\in D_4$ satisfies it, proving equality of
the two intersections and the displayed inverse image.

Finally
$h_0(\sqrt2(a,a,a),\sqrt2(b,b,b))=6h(a,b)$.
The four $\delta_j$ form a basis of $D_4$, since their determinant
has absolute value two and $D_4$ has index two in $\mathbb Z^4$.
Their scalar products give \eqref{wcl:eq:common-gram}.
The least nonzero norm in $D_4$ is two: a norm-one integer vector
has odd coordinate sum, while $(1,-1,0,0)$ has norm two.
This proves the exact minimum and determinant.
\end{proof}

The two slice indices are also literal. Define the common
four-dimensional real subspace
$\mathcal U=\{(u,u,u):u\in U\}$. The proof gives
\[
 \mathcal L(N)\cap\mathcal U
 =\{(k,k,k)/\sqrt2:k\in P\},\quad
 P=2\mathbb Z^4+\mathbb Z(1,1,1,1),
\]
\[
 \mathcal L(\Lambda)\cap\mathcal U
 =\{(k,k,k)/\sqrt2:k\in P_0\},\quad
 P_0=\{k\in P:\textstyle\sum k_i\equiv0\pmod4\}
     =2D_4+\mathbb Z(1,1,1,1).
\]
The first proof does not need Wilson membership to obtain these
slice equations. The lattice $P$ has basis
$2e_0,2e_1,2e_2,(1,1,1,1)$ and covolume $8$.
The lattice $2D_4$ has covolume $2^4\cdot2=32$, so
$[P:2D_4]=4$. The class of $(1,1,1,1)$ has exact order two
modulo $2D_4$, yielding $[P_0:2D_4]=2$. Therefore
\[
 [\mathcal L(N)\cap\mathcal U:\mathcal I]=4,\qquad
 [\mathcal L(\Lambda)\cap\mathcal U:\mathcal I]=2.
\]
The rank-four intersection, strictly below rank $24$, also proves
that these full lattices are not commensurable \emph{in the retained
embedding}: a common finite-index subgroup would have full rank.
Their exact existing relation is the common lattice with the
explicit embeddings, inverse map, metrics, and indices just proved.

\subsection{The cubic value group on the literal common lattice}
\label{wcl:sec:common-cubic}

Retain $F(y)=\tau(y_1,y_2,y_3)$, so the Albert off-diagonal
determinant is $2F$. For $y=\sqrt2(a,a,a)\in\mathcal I$ write
$a=a_0+\mathbf a$ in the same quaternionic $U$. Since
$\mathbf a^2=-N(\mathbf a)$ and $a_0$ is real, direct multiplication
inside this associative subalgebra gives
\[
 \operatorname{Re}(a^3)
 =a_0^3-3a_0N(\mathbf a)=4a_0^3-3a_0N(a).
\]
Every scalar and factor is therefore
\begin{equation}
 F(\sqrt2(a,a,a))
 =2\sqrt2\bigl(4a_0^3-3a_0N(a)\bigr).
 \label{wcl:eq:common-cubic}
\end{equation}
For $a\in D_4$, $N(a)$ is even, because
$N(a)\equiv\sum a_i=0\pmod2$. The integer in parentheses is
therefore even, so every value belongs to $4\sqrt2\mathbb Z$.
For the actual lattice vector $a=(1,1,0,0)$ it equals
$4-3\cdot2=-2$, giving $F=-4\sqrt2$.
Consequently the \emph{additive subgroup generated by all values}
is exactly
\[
 \bigl\langle F(\mathcal I)\bigr\rangle_{\mathbb Z}
       =4\sqrt2\mathbb Z,\qquad
 \bigl\langle d(Q_{\mathbb O}(0;\mathcal I))\bigr\rangle_{\mathbb Z}
       =8\sqrt2\mathbb Z.
\]
This assertion concerns the generated additive groups; it does not
claim that every element of either group is an individually attained
cubic value.
For an explicit distinction, put
$P(a)=a_0(a_0^2-3(a_1^2+a_2^2+a_3^2))$.
If $3\mid P(a)$, reduction modulo three forces $3\mid a_0$,
and substitution $a_0=3b$ gives
$P(a)=9b(3b^2-a_1^2-a_2^2-a_3^2)$.
Thus $3\mid P(a)$ implies $9\mid P(a)$.
The value $P(a)=6$ is impossible, so $F=12\sqrt2$ is not
attained, although it belongs to the proved additive value group.


\subsection{Primary source and exact reproducibility data}

Wilson's octonionic construction is the source of the pair and total
conditions \cite[Sections~2--4, pp.~2--6]{esn:Wilson2009}. His page~2
uses the odd-minus lattice $L$ and $R=\bar L$; page~3 gives pair sums
in $L\bar s$, the total in $Ls$, and the metric equal to one half of
the sum of octonion norms. The displayed basis map and
\eqref{wcl:eq:metric-isometry} retain these conjugations and the exact
scale in the present coordinates. The distinct coordinate change on
Wilson's page~6 is not substituted for this isometry.

The complete integer-arithmetic program and its output are supplied in the
following directory (the two displayed lines form one path):
\begin{quote}\small
\path{supporting_materials/workbench/evidence/}\\
\path{lattice_foundation_20260906/wilson/}
\end{quote}
The files are \path{check_wilson.py} and \path{EXACT_CHECKS.json}.
They record the generating-lattice index $4096$, Gram determinant $1$,
even diagonal, integral cyclic action, the fixed and coinvariant Gram
matrices, and the index $6561$ of their sum. The same program checks
all $64$ multiplication-table products for
\eqref{wcl:eq:actual-basis-map} and the stated exact cubic tests.
The proofs above establish rootlessness and the global descriptions;
these conclusions are not inferred from a finite search.
